\chapter{Introduzione}

\section{Statement (bozza iniziale)}
Il progetto propone un'applicazione locale per esplorare e organizzare un catalogo multimediale. Il suo interesse principale, nel contesto dell'ingegneria del software, è la progettazione di un nucleo riutilizzabile per cataloghi e la sua applicazione concreta a un catalogo di film.

Il nucleo dovrà trattare in modo generico profili utente, preferenze per media e autori, sezioni, media riferiti a file generici e autori o contributori associati alle opere. La specializzazione concreta rappresenterà film, generi, file video e locandine. In questo modo le scelte specifiche del cinema non dovranno coincidere con le responsabilità del framework riutilizzabile.

\section{Obiettivo e confini}
Le funzioni approvate comprendono gestione dei profili e delle preferenze, ricerca e navigazione del catalogo, inserimento di film, generi e autori e le relative associazioni. Le schermate consentono di esercitare queste funzioni, ma costituiscono un supporto al progetto: l'attenzione è rivolta al modello di dominio, ai servizi applicativi, alla persistenza e alla verifica.

Il sistema è pensato per un uso locale e per file selezionati dal computer dell'utente. La riproduzione dei video non è stata definita come requisito. Sono ancora da concordare, tra gli altri dettagli, la gestione fisica dei file, alcune cardinalità e il comportamento delle cancellazioni; la relazione non le assume in anticipo.

\section{Organizzazione della relazione}
La relazione documenterà contesto e requisiti, casi d'uso e disegno funzionale, architettura e modelli dei dati, dettagli implementativi significativi e verifiche funzionali e strutturali. Diagrammi, template e risultati saranno inseriti quando corrispondono a decisioni progettuali e prove effettivamente svolte.

\section{Dichiarazione degli strumenti di supporto}
Il disciplinare richiede di documentare gli strumenti di intelligenza artificiale generativa eventualmente impiegati, le modalità d'uso e le parti del lavoro interessate. Questa sezione dovrà essere completata dall'autore in base all'uso effettivo degli strumenti e rivista prima della consegna.
